\documentclass[11pt]{article}
\usepackage[a4paper,margin=1in]{geometry}
\usepackage{amsmath, amssymb, amsthm}
\usepackage{bm}
\usepackage{hyperref}
\hypersetup{colorlinks=true, linkcolor=blue!60!black, urlcolor=blue!60!black}

\newcommand{\ev}[1]{\left\langle #1 \right\rangle}
\newcommand{\given}{\,\middle|\,}
\newcommand{\dd}{\mathrm{d}}
\newcommand{\Tr}{\mathrm{Tr}\,}
\newcommand{\T}{\mathsf{T}}
\newcommand{\Tinv}{\mathsf{T}^{-1}}
\newcommand{\Sig}{\Sigma_V}
\newcommand{\Sinv}{\Sig^{-1}}

\title{Laplace-tilted importance sampling for finite-$T$ cluster spin DMFT}
\author{Notes on the implementation in \texttt{CspinDMFT\_finite\_T}}
\date{}

\begin{document}
\maketitle

\section{Setting and the failure mode}

The self-consistent Monte Carlo loop draws a mean-field noise trajectory
$V = \{V_i^\alpha(\tau)\}_{i,\alpha,\tau}$ from a multivariate Gaussian
\begin{equation}
  p(V) \;=\; \mathcal{N}\!\bigl(\mu,\;\Sig\bigr),
  \label{eq:p}
\end{equation}
where $\mu$ is the (static) mean field and $\Sig$ is the mean-field
covariance built from the previous iteration's spin correlations
(the Frequency-domain block-diagonal structure already implemented by
\texttt{FrequencyCovarianceCluster}). For each sampled $V$ we compute
the imaginary-time-ordered propagator and its trace
\begin{equation}
  Z(V) \;=\; \Tr\, T_\tau \exp\!\Bigl(-\!\!\int_0^\beta H[V(\tau)]\,\dd\tau\Bigr),
  \qquad
  H[V] = J\,\mathbf{S}\!\cdot\!\mathbf{S} + V\!\cdot\!\mathbf{S}.
  \label{eq:Z}
\end{equation}
The observable of interest is
\begin{equation}
  \ev{\mathcal{O}} \;=\;
  \frac{\mathbb{E}_p\!\left[\,Z(V)\,\mathcal{O}(V)\,\right]}{\mathbb{E}_p\!\left[\,Z(V)\,\right]},
  \qquad
  \mathcal{O}(V) \equiv \frac{\Tr\!\left[F(V)\,S\,B(V)\,S\right]}{Z(V)},
  \label{eq:estimator}
\end{equation}
estimated by the self-normalised Monte Carlo ratio
$\widehat{\ev{\mathcal{O}}} = \sum_s \mathrm{Tr}[FSBS]_s\,/\sum_s Z_s$
with $V_s \sim p$.

\paragraph{The effective sample size collapses with $\beta$.}
The variance of the estimator is controlled by the spread of the
importance weights $w_s = Z(V_s)$ under $p$. A useful diagnostic is the
effective sample size
\begin{equation}
  \mathrm{ESS} \;\equiv\; \frac{\bigl(\sum_s w_s\bigr)^2}{\sum_s w_s^2} \;\in\;[1,M].
  \label{eq:ess}
\end{equation}
For a single spin in field $V$ with $Z(V)=2\cosh(\beta\!\left|V\right|/2)$
and $V\!\sim\!\mathcal{N}(0,\sigma^2)$,
\begin{equation}
  \frac{\mathrm{Var}_p[Z]}{\mathbb{E}_p[Z]^2}
  \;\approx\; \exp\!\Bigl(\tfrac{\beta^2\sigma^2}{4}\Bigr) - 1,
  \label{eq:varZ}
\end{equation}
so $\mathrm{ESS}/M \sim \exp(-\beta^2\sigma^2/4)$.
For the cluster the exponent picks up an additive factor in the number of sites,
so $\mathrm{ESS}/M$ degrades \emph{exponentially} in both $\beta$ and $N$.
This is the origin of the ``huge stds at low $T$'' failure mode and of
the overflow NaNs observed in the \texttt{--uncoupled\_spins} branch.

\section{The ideal proposal and why it is intractable}

For any positive proposal $q(V) > 0$,
\begin{equation}
  \ev{\mathcal{O}}
  \;=\;
  \frac{\mathbb{E}_q\!\left[\,w(V)\,Z(V)\,\mathcal{O}(V)\,\right]}
       {\mathbb{E}_q\!\left[\,w(V)\,Z(V)\,\right]},
  \qquad
  w(V) \equiv \frac{p(V)}{q(V)}.
  \label{eq:reweighted}
\end{equation}
The zero-variance choice is $q^{*}(V) \propto p(V)\,Z(V)$, for which
$w(V)\,Z(V) \equiv \mathrm{const}$ and the estimator becomes the plain
mean of $\mathcal{O}$ over $V_s\sim q^*$. We cannot sample from $q^*$
directly because $Z(V)$ is the partition function of a non-trivial
imaginary-time-ordered product, with no closed form.

\section{Laplace approximation to \texorpdfstring{$q^*$}{q*}}

We expand $\log p(V) + \log Z(V)$ to second order around its joint
maximum $V_\star$, giving the Gaussian proposal
\begin{equation}
  q(V) \;=\; \mathcal{N}\!\bigl(V_\star,\, H^{-1}\bigr),
  \qquad
  H \;=\; \Sinv \;-\; \chi_V[V_\star],
  \label{eq:H}
\end{equation}
where
\begin{align}
  V_\star \;\text{satisfies}\;\;& \Sinv\,V_\star \;=\; \ev{\mathbf{S}}_{V_\star},
  \label{eq:saddle}\\
  \chi_V[V](\tau,\tau')_{i\alpha,j\beta}
  \;\equiv\; & \frac{\partial^{2}\log Z[V]}{\partial V_i^\alpha(\tau)\,\partial V_j^\beta(\tau')}.
  \label{eq:chiV}
\end{align}
The DMFT self-consistency \emph{already} solves \eqref{eq:saddle} for the
mean, so $V_\star = \mu$ for free. Only the Gaussian \emph{covariance}
changes: it becomes wider than $\Sig$ by the susceptibility correction.

\subsection{Continuum and discrete forms of $\chi_V$}
For a static field $V$, $\partial \log Z/\partial V_i^\alpha = -\beta \ev{S_i^\alpha}_V$.
For an imaginary-time-dependent field one finds, in the continuous
formulation,
\begin{equation}
  \frac{\partial^{2}\log Z[V]}{\partial V_i^\alpha(\tau)\,\partial V_j^\beta(\tau')}
  \;=\; \chi_S(\tau,\tau')_{i\alpha,j\beta},
  \qquad
  \chi_S \;\equiv\; \ev{T_\tau S_i^\alpha(\tau)\,S_j^\beta(\tau')}_{\!\rm conn}.
  \label{eq:chiS}
\end{equation}
On the discrete time grid $\tau_n = n\,\Delta\tau$, $\Delta\tau = \beta/N_t$,
where $\partial V_i^\alpha(\tau_n)$ acts as if integrated against
$\Delta\tau$,
\begin{equation}
  \chi_V \big|_{\text{disc}} \;=\; (\Delta\tau)^{2}\,\chi_S,
  \label{eq:dt2}
\end{equation}
which gives the implementation factor \verb!tilt_scale = dt*dt!.
With \verb!tilt_scale = 1! the quadratic form in the weight overflows
within the first sample; with \verb!tilt_scale!~$=\Delta\tau^{2}$ the
ESS plateau increases monotonically (\S\ref{sec:results}).

\subsection{Block-wise Matsubara implementation}
$\Sig$ and $\chi_V$ are translationally invariant in imaginary time,
hence block-diagonal in the cosine/sine Matsubara basis used by
\texttt{FrequencyCovarianceCluster}. For each frequency $n$ and symmetry
block $b$,
\begin{equation}
  H_n^{(b)} \;=\; \bigl[\Sig\,_n^{(b)}\bigr]^{-1}
              \;-\; \texttt{tilt\_scale}\;\chi_S\,_n^{(b)},
  \qquad
  \Sig\,_n^{\rm tilted (b)} \;=\; \bigl[H_n^{(b)}\bigr]^{-1}.
  \label{eq:Hblock}
\end{equation}
If any block fails the positive-definiteness check
$\lambda_{\min}(H_n^{(b)}) \le 0$ we keep the untilted block for that
frequency. This robust fallback matters near a phase transition where
$\chi$ can become large enough to invert the sign of an eigenvalue.
The cosine-Fourier convention used to build $\chi_n$ matches the one
used by \texttt{fill()}, so the quadratic form below is consistent
without an extra normalisation factor.

\section{Per-sample weight}

Because $p$ and $q$ share the mean $\mu = V_\star$,
\begin{align}
  \log w(V) \;&=\; \log p(V) - \log q(V) \nonumber\\
  \;&=\; -\tfrac{1}{2}\,(V-\mu)^{\T}\bigl[\Sinv - H\bigr](V-\mu)
        \,-\, \tfrac{1}{2}\log\det\!\bigl(\Sig\,H\bigr) \nonumber\\
  \;&=\; -\tfrac{1}{2}\,\texttt{tilt\_scale}\,(V-\mu)^{\T}\chi_S(V-\mu)
        \;+\; C,
  \label{eq:logw}
\end{align}
with $C$ a sample-independent constant that drops out of any
self-normalised estimator.
In the orthonormal cosine/sine basis used to draw $V$, the quadratic
form decouples across Matsubara frequencies,
\begin{equation}
  (V-\mu)^{\T}\chi_S(V-\mu)
  \;=\; \sum_{n} \tilde V_n^{\T}\,\chi_S\,_n\,\tilde V_n,
  \label{eq:quad}
\end{equation}
which is what \verb!compute_log_pq_weight()! evaluates after a forward
Fourier of the sampled time-domain trajectory.

\paragraph{Why the sign is correct.}
Tail samples of $V$ have anomalously large $Z(V)$. The optimal $q^*$
puts more probability there, so when we draw from $q$ those samples are
oversampled relative to the target $p$. The compensating weight
$w = p/q$ must therefore be \emph{smaller} in the tail — exactly the
behaviour of \eqref{eq:logw} for \verb!tilt_scale!~$>0$.

\section{What enters the running accumulators}

Every per-sample quantity that the naïve code multiplies by $Z$ gets
replaced by the product $w\,Z$, and every per-sample quantity multiplied
by an unnormalised trace $x = \Tr[FSBS]$ by $w\,x$:
\begin{align}
  \text{partition function:}\quad &
  \sum_s Z_s \;\longrightarrow\; \sum_s w_s\,Z_s, \\
  \text{correlator estimator:}\quad &
  \sum_s x_s \;\longrightarrow\; \sum_s w_s\,x_s, \\
  \text{Delta-method second moments:}\quad &
  \sum_s x_s^{2},\ x_s Z_s,\ Z_s^{2}
  \;\longrightarrow\;
  \sum_s (w_s x_s)^{2},\ (w_s x_s)(w_s Z_s),\ (w_s Z_s)^{2},
\end{align}
which preserves the Delta-method formula for the std of the ratio
estimator and the ESS formula \eqref{eq:ess} with $w_s\,Z_s$ in place of
$Z_s$. In the implementation a single argument \verb!weight! to
\verb!compute_spin_observables()! applies these substitutions
consistently across all accumulators.

\section{Diagnostics and convergence}
\label{sec:results}

For the $N=2$, $J=0.5$ square cluster at $M = 4\times10^4$ samples
per DMFT iteration, the converged ESS plateaus are
\begin{center}
\begin{tabular}{c|cc|c}
$\beta$ & naïve $\mathrm{ESS}/M$ & tilted $\mathrm{ESS}/M$ & variance reduction \\
\hline
$1$ & $0.994$ & $1.000$ & $\sim$1$\times$ \\
$3$ & $0.751$ & $0.947$ & $\sim$1.3$\times$ \\
$5$ & $0.341$ & $0.894$ & $\sim$2.6$\times$ \\
\end{tabular}
\end{center}
Mean correlators agree within sample stds at all $\tau$, all $\beta$
(\verb!Plots/tilted_vs_naive_N2_J0.5.{pdf,png}!), confirming that
reweighting is unbiased and that the variance reduction is genuine.

\section{Open issues}

\begin{enumerate}
  \item \textbf{Susceptibility from spin correlations.} The
        implementation currently builds $\chi_S$ from the previous
        iteration's \verb!my_spin_correlations_Re! with the disconnected
        $\ev{S}\ev{S}$ subtracted. This is the cluster susceptibility
        of the \emph{bath-dressed} impurity, not the susceptibility at
        the static saddle $V_\star = \mu$. They agree at
        self-consistency and differ away from it, which is why
        ESS$/M$ at $\beta=5$ plateaus near $0.9$ rather than $1.0$. A
        deterministic propagator run at zero noise would yield the
        Laplace-exact $\chi$.
  \item \textbf{Reduced layout for $\chi$.} The
        \verb!correlation_categories! layout stores only representative
        pairs; \verb!FrequencyCovarianceCluster::fill! requires an
        all-pairs CluCorrTen. The current
        \verb!build_connected_susceptibility! reconstructs missing
        pairs by translation fallback. For non-translationally-invariant
        problems (open boundaries, disorder) a proper mapping is
        required.
  \item \textbf{Positive definiteness near criticality.} The block-wise
        PD check guards against $\chi_S\,_n^{(b)}$ exceeding
        $[\Sig\,_n^{(b)}]^{-1}$ in any direction. At a phase transition
        this guard activates for the dominant mode, the tilt is
        suppressed there, and ESS recovery degrades smoothly rather
        than blowing up. A monitoring counter (number of untilted
        blocks per iteration) is worth promoting to a stored HDF5
        diagnostic.
  \item \textbf{Uncoupled branch.} \texttt{--tilted\_is} is currently
        gated on \verb!!my_pspace.uncoupled_spins!. Extending it to the
        uncoupled branch is straightforward (the per-site susceptibility
        is analytic) and would address the overflow NaN at
        $\beta\!\gtrsim\!5$ at its statistical source rather than via
        log-space bookkeeping.
  \item \textbf{Compounding with cluster size.} The relative-variance
        bound \eqref{eq:varZ} adds in the number of sites. The
        $N=2$ improvement from $0.34$ to $0.89$ at $\beta=5$ would
        compound to roughly $\bigl(0.89/0.34\bigr)^{\,N/2}$ across
        $N$ sites if the tilt remains effective per site. This is
        what makes the Laplace tilt structurally important for
        production-scale $N=9$ runs.
\end{enumerate}

\end{document}
